\documentclass[11pt]{article}
\usepackage[T1]{fontenc}
\usepackage[utf8]{inputenc}
\usepackage{lmodern,amsmath,amssymb,microtype}
\usepackage[letterpaper,margin=0.87in]{geometry}
\usepackage[colorlinks=true,urlcolor=blue]{hyperref}
\hypersetup{pdftitle={Clique-separator weight, layered tree-independence, and fractional fragility}}
\setlength{\parindent}{0pt}\setlength{\parskip}{5pt}
\newcommand{\done}{\hfill\(\square\)}
\begin{document}
\begin{center}{\Large\bfseries Clique-separator weight, layered tree-independence,\\ and fractional fragility}\par
\small Research note \enspace$\cdot$\enspace 18 November 2026\end{center}

Galby, Munaro and Yang ask whether bounded layered tree-independence guarantees balanced clique-based separators of sublinear weight. We obtain a strongly sublinear bound for each fixed layered width. We also answer two questions about converse implications using a hereditary class derived from Moshkovitz and Shapira's graphs. [\href{https://arxiv.org/abs/2402.18352}{Galby, Munaro and Yang, Questions 3 and 4}; \href{https://arxiv.org/abs/2506.12424}{Dallard et al., Questions 1.2 and 7.2}]

A clique-based separator is a family of disjoint cliques whose removal splits the remaining vertices into two anticomplete sets of size at most $2n/3$ each. Its weight is $\sum_C\log_2(|C|+1)$. A layering partitions the vertices into indexed layers, with all edges within one layer or between consecutive layers. Layered tree-independence at most $k$ means there is a tree decomposition and a layering such that every bag meets every layer in a set of independence number at most $k$.

\textbf{Theorem 1.} Every $n$-vertex graph of layered tree-independence at most $k\ge1$ admits a clique-based separator of weight
\[
                    O_k\big(n^{k/(k+1)}\log(n+1)\big).
\]
In particular Galby, Munaro and Yang's Question 4, restated as Question 1.2 by Dallard et al., has an affirmative answer.

\textit{Proof.} A graph on $s\ge1$ vertices of independence number at most $k$ partitions into $O_k(s^{1-1/k})$ cliques. For $k=1$ it is complete. For $k\ge2$, the classical Ramsey bound $R(k+1,p)\le\binom{k+p-1}{k}$ guarantees a clique of size at least $c_k v^{1/k}$ among any $v$ remaining vertices for a suitable positive $c_k$. Repeatedly remove such a clique. While the number remaining lies in $(u/2,u]$, at most $O_k(u^{1-1/k})$ removals are possible. Sum these bounds over $u=s,s/2,s/4,\ldots$.

Choose $t=\lceil n^{1/(k+1)}\rceil$ and delete the layers of a smallest residue class of indices modulo $t$. Their vertices $Z$ cost at most $n/t$ in singleton cliques. Every component left occupies fewer than $t$ consecutive layers. If a component $H$ has more than $2n/3$ vertices it is unique. Restrict the tree decomposition to $H$ and take a bag $B$ whose removal leaves components of size at most $|H|/2$. For the bag lemma, assign each vertex of $H$ to one node containing it and take a weighted centroid of the decomposition tree. The bags of a component avoiding the centroid lie in a single branch, by the connected-subtree and edge axioms. The centroid bag thus has the required property. If no component is heavy, take $B=\varnothing$.

Every component outside $Z\cup B$ now has size at most $2n/3$. Such components can be grouped into two sets of size at most $2n/3$: take a component of size at least $n/3$ alone if one exists. Otherwise accumulate components to reach $n/3$ without exceeding $2n/3$, or use all of them if the total is smaller.

Split $B\subseteq H$ by its layers. If the nonempty intersection sizes are $s_1,\ldots,s_a$, then $a\le t$, $\sum s_i\le n$, and each intersection has independence number at most $k$. Partitioning each separately costs at most
\[
 O_k\!\left(\log(n+1)\sum_{i=1}^a s_i^{1-1/k}\right)
       \le O_k\!\left(\log(n+1)t^{1/k}n^{1-1/k}\right).
\]
For $k=1$ the sum counts the nonempty intersections. Adding $n/t$ gives the claimed bound.\done

\textbf{Corollary.} Every fractionally tree-$\alpha$-fragile graph class admits clique-based separators of weight $o(n)$.

\textit{Proof.} For every fixed integer $r\ge2$, fragility supplies a distribution on deletion sets with inclusion probability at most $1/r$ per vertex, such that every remaining graph has a tree decomposition with bag independence at most a class-dependent $f(r)$. Take a deletion set of size at most $n/r$ by averaging. A centroid bag in the remainder separates it into components of size at most $n/2$. Its independence number is at most $f(r)$. The Ramsey partition from the proof above covers this bag with weight $o_r(n)$ for fixed $r$. The singleton deletion set costs at most $n/r$. The worst-case weight divided by $n$ thus has limit superior at most $1/r$ for every $r$.\done

\textbf{Theorem 2.} There is a hereditary triangle-free graph class with clique-based separator weight $o(n)$, unbounded independence degeneracy, and no fractional tree-$\alpha$-fragility. This answers Dallard et al.'s Question 7.2 and Galby, Munaro and Yang's Question 3 negatively.

\textit{Proof.} Moshkovitz and Shapira's Theorem 1 provides graphs $G_{N,d}$ for $2\le d\le(\log_2\log_2 N)/648$ with average degree at least $d$ and girth at least $L=(\log_2 N)/(6d)^2$. Every subgraph on $s\ge L$ vertices has a $2/3$-balanced edge separator of size at most $s(\log_2\log_2 s)^2/\log_2 s$. Use $d\to\infty$ and choose $N$ so large for each $d$ that $L>4$. Take all induced subgraphs of these graphs. [\href{https://arxiv.org/abs/1502.00409}{Moshkovitz and Shapira, Theorem 1}]

Every member is triangle-free. An $s$-vertex subgraph with $s<L$ is a forest and has a balanced vertex separator of size at most one, by a tree centroid in its heavy component. For $s\ge L$, choose one endpoint of each separator edge. The resulting vertex set separates the two parts and incurs at most the edge-separator bound. Singletons have weight one, proving the uniform $o(s)$ bound. It also bounds the independence number of the separator, as required in Question 3.

Every graph of average degree at least $d$ has a nonempty induced subgraph of minimum degree at least $d/2$, by repeatedly deleting vertices of lower degree. All neighborhoods in this subgraph are independent because the graph is triangle-free. Thus its minimum neighborhood independence number is at least $d/2$. This is a lower bound on the independence degeneracy, which takes the maximum of that quantity over induced subgraphs.

Finally suppose that the class were fractionally tree-$\alpha$-fragile and fix $r=4$. In any remaining graph all bags have at most $R(3,f(4)+1)-1$ vertices, by triangle-freeness and Ramsey's theorem. Thus its treewidth, and hence its degeneracy, is bounded by a class constant $M$. It has at most $MN$ edges. Each edge of $G_{N,d}$ survives with probability at least $1-2/4=1/2$, so the expected number of surviving edges is at least $dN/4$. Taking $d>4M$ is a contradiction.\done

\small
\textbf{Scope and status.} The positive argument combines sparse layer deletion with separate Ramsey partitions of a single bag's layer intersections. Dallard et al., Proposition 5.4(c), already apply Ramsey's theorem to a bag of bounded independence number. The negative construction itself is due to Moshkovitz and Shapira; the deduction here concerns the two later questions. The latter counterexample supplies a sublinear bound of logarithmic relative order. Dvo\v r\'ak's question with a strongly sublinear hypothesis has a different assumption. Searches through 18 November 2026 found no previous answer to these three questions in the stated forms. The arguments await independent review. [\href{https://arxiv.org/abs/2506.12424}{Dallard et al.}; \href{https://arxiv.org/abs/2402.18352}{Galby, Munaro and Yang}]
\end{document}
